% !TEX program = xelatex
\documentclass[11pt,a4paper]{article}

\usepackage[a4paper,margin=22mm]{geometry}
\usepackage{xeCJK}
\usepackage{fontspec}
\usepackage{amsmath,amssymb}
\usepackage{booktabs}
\usepackage{tabularx}
\usepackage{array}
\usepackage{enumitem}
\usepackage{xcolor}
\usepackage{pgfplots}
\usepackage{hyperref}

\pgfplotsset{compat=1.18}

\setCJKmainfont{Hiragino Mincho ProN}
\setCJKsansfont{Hiragino Sans}
\setCJKmonofont{Hiragino Sans}
\setmonofont{Menlo}[Scale=0.82]

\xeCJKDeclareCharClass{CJK}{"2460 -> "24FF, "2013 -> "2015, "2018 -> "201F, "2026, "2500}

\hypersetup{colorlinks=true, linkcolor=black, urlcolor=blue!60!black,
  pdftitle={案E 第一次結果 —— 齢を状態に持つ方策は、模倣が届かない遅延領域を取れるか},
  pdfauthor={児玉靖司}}

\newcommand{\code}[1]{\texttt{#1}}
\newcommand{\prob}[1]{\textbf{\textcolor{red!55!black}{#1}}}
\newcommand{\keyp}[1]{\textbf{\textcolor{blue!45!black}{#1}}}

\title{\vspace{-12mm}
  \Large 案E 第一次結果\\[2mm]
  \large 齢を状態に持つ方策は、模倣が届かない遅延領域を取れるか\\[1mm]
  \normalsize ―― 3つの実験、そこで見つかった実験設計の誤り、および実装の誤り\\[1mm]}
\author{\normalsize 児玉靖司}
\date{\normalsize 2026年8月14日}

\begin{document}
\maketitle
\vspace{-6mm}

\begin{abstract}
\noindent
「機械学習・強化学習によるロボットアーム／フィジカル AI 研究」（2026年8月12日）で示した
\keyp{案E}に着手し、\keyp{機材を買う前にシミュレータだけで出せる一次結果}をここに報告する。
1 関節・1\,kHz の位置サーボに対し、上位からの目標が遅れて届く状況を作り、
\keyp{模倣学習と強化学習 $\times$ 齢の有無}の4条件を\emph{同一容量・同一軌道}で比較した。

結果は\keyp{仮説を半分支持し、半分否定した}。
支持された側 ―― \keyp{遅延 400\,ms 以上では強化学習が模倣学習を 5 種すべてで上回る}
（400\,ms で $-28\,\%$、800\,ms で $-15\,\%$）。
案Eの中心にある「\emph{教師が失敗する領域は模倣では学べず、そこは強化学習が取る}」が
実際に観測された。
否定された側 ―― \prob{齢の効果は遅延が大きいほど薄れ、400\,ms では逆に害になった}
（模倣で $+22\,\%$、勝ち 1/5 種）。

否定の原因は方策側ではなく\prob{実験設計}にあった。
遅延を1回の走行のあいだ一定にしたため、\keyp{齢は「定数＋配布周期の位相」に退化し、
定数部分は学習で重みに吸収される}。案Eの主張は「\emph{遅延が変動する条件下で}」であり、
掃引すべき軸は遅延の大きさではなく\keyp{遅延のばらつき}であった。

そこで軸を替えて実験2を行い、\keyp{主張は支持された}。
文脈生成のジッタを $\pm 0$ から $\pm 200$\,ms へ広げると、齢の効果は
\keyp{$0\,\%$ から $-31\,\%$ へ単調に強まり、$\pm 50$\,ms 以上では5種すべてで齢ありが勝つ}。
すなわち\keyp{齢が効くのは遅延が大きいときではなく、遅延がばらつくときである}。
さらに実験3で $a$ と $\hat a$ を分離した ―― \keyp{主役は齢 $a$}（ジッタ $\pm 200$\,ms で $-28\,\%$）、
\keyp{$\hat a$ 単独の寄与は小さいが符号は一貫して正}（$-6\,\%$、5/5 種）。
効果が最大になるのは\keyp{文脈の脱落}（メールボックス溢れ）を入れたときで、
脱落 30\,\% では\keyp{$-49\,\%$（5/5 種）}に達する。
一方\prob{強化学習側では3実験とも齢の効果が検出できず}、
CEM が齢を使う解に届いていない疑いが残った。

途中で\prob{実装の誤り}も1つ見つかった。$\hat a$ を「次に本当に届く時刻」から作っていたため、
\prob{OS には出せない情報（起動区間の目印）が漏れていた}。
名目周期から作る正しい実装に直すと、ジッタ 0 での「齢の効果 $-10\,\%$」は消えた。
\keyp{OS が何を答えられるかを厳密にしないと、齢の効果を過大評価する}という、
案E そのものに関わる教訓である。
\end{abstract}

\vspace{-2mm}
{\small\tableofcontents}

\section{これは何か}

\code{docs/ml\_rt\_proposals.pdf} 第8節（案E）の着手ぶんである。
案Eは3年・380万円の計画だが、\keyp{その一次結果は装置を買わずに出せる}と書いた。
本稿はそれを実行した記録であり、実装は \code{xinu-mesh-rtos/planE/}（純 Python、依存ゼロ）にある。

問いは次の一文である。

\begin{quote}
\keyp{遅延を「情報として与えられた」方策は、遅延を「上界としてのみ扱う」方策より
どれだけ良くなれるか。そして模倣が原理的に届かなくなる領域を、強化学習は取れるか。}
\end{quote}

案D（\code{vla\_rt\_review.pdf} 第III部）は「齢が効く」ことを特権教師つき模倣学習で示す計画だが、
\prob{教師自身が失敗する大遅延領域は模倣では学べない}と自ら範囲外を宣言している。
本実験はその境界を実測し、その先を強化学習が取れるかを見る。

\section{実験系}

\subsection{6つの条件}

\begin{table}[htbp]
\centering
\small
\caption{比較する6条件。方策はすべて同一構造・同一パラメータ数（121）}
\begin{tabularx}{\textwidth}{@{}l X l@{}}
\toprule
\textbf{条件} & \textbf{制御} & \textbf{齢 $a$ と $\hat a$} \\
\midrule
\code{teacher} & 遅延ゼロの理想制御（先行補償つき PD）。\keyp{性能の下界} & 本当の目標を知っている \\
\code{naive}   & 最後に知った目標角をそのまま指令。\keyp{学習なしの基準線} & 使わない \\
\code{P0} & 模倣学習（特権教師つき、DAgger 1 巡） & なし（端子はゼロ） \\
\code{P1} & 模倣学習 & \keyp{あり} \\
\code{P2} & 強化学習（CEM） & なし（端子はゼロ） \\
\code{P3} & 強化学習（CEM） & \keyp{あり} \\
\bottomrule
\end{tabularx}
\end{table}

\noindent
方策は入力 8 $\to$ 隠れ 12(tanh) $\to$ 出力 1 の\keyp{121 パラメータ}の残差方策で、
基底（＝最後に知った目標角）の上に $\pm 40$\,deg の差分を出す。
1\,kHz で回すこと、将来ノード上で固定小数点にすることを見込んだ規模である。

\subsection{案Eに固有の入力 $\hat a$}

案Dが渡すのは「いま何 µs 古いか」（齢 $a$）だけである。本実験はもう1つ、
\keyp{「あと何 µs で新しい文脈が来るか」}$\hat a$ を渡す。
tickless ワンショットタイマは次の起床時刻を持っているので、
\keyp{OS だけが原理的に答えられる値}である。方策から見れば未来の情報であり、
「待つ」という行動を選べるかがこの1次元で変わる。
DA-MDP・世界モデル・遅延ロバスト化のいずれにもない入力である。

\subsection{環境}

\begin{itemize}[leftmargin=1.4em,itemsep=2pt]
  \item サーボ：位置指令 $\to$ 内部 PD（$\omega=30$\,rad/s, $\zeta=1.0$）、
        角速度飽和 \code{VMAX}=180\,deg/s、可動域 0--180\,deg、制御周期 1\,ms。
  \item 上位目標：0.15--0.8\,Hz の滑らかな軌道。サーボ帯域（4.8\,Hz）より十分低いので、
        \keyp{遅延ゼロなら誤差ほぼ 0 で追える}＝教師が「正解」として機能する。
  \item 文脈配布：周期 100\,ms（10\,Hz）で生成され、輸送遅延 $d$ を経て届く。
        齢は $d$ から $d+100$\,ms を鋸歯状に動く。
  \item \keyp{実物の VLA は使わない}（案Dの発想 ―― 遅延の研究に、遅延を生む実物は要らない）。
\end{itemize}

\subsection{しきい値を先に宣言する}

後から都合よく選ばないため、\code{sweep.py} の冒頭に固定してある。

\begin{itemize}[leftmargin=1.4em,itemsep=2pt]
  \item \code{TOL}\,$=2.0$\,deg（可動範囲 50\,deg p-p の 4\,\%）
  \item $d^{*}_{IL}$ ＝ P1 の rms が \code{TOL} を超える最小の遅延
  \item $d^{*}_{RL}$ ＝ P3 の rms 中央値が P1 の rms 中央値を下回る最小の遅延
\end{itemize}

\section{実験系そのものの検証}

測っているものが「齢の効果」でなくなる壊れ方は3つある。
\code{test\_planE.py} が機械的に検査する（17 項目、\keyp{全通過}）。

\begin{enumerate}[leftmargin=1.8em,itemsep=3pt]
  \item \prob{容量が揃っていない} $\to$ 測っているのは容量差。
        4条件のパラメータ数が同一（121）であることを検査する。
        \keyp{齢なし条件も端子は持ったままゼロを流す}。
  \item \prob{齢マスクが効いていない} $\to$ 齢なし条件が齢を見てしまう。
        齢を変えても入力が不変であること、齢あり条件では出力が実際に変わること
        （$\Delta=-6.85$\,deg）、齢の端子以外は同一であることを検査する。
  \item \prob{学習と評価が同じ軌道} $\to$ 測っているのは過学習。
        学習は種 2000 系列、評価は 1000 系列で分離する。
\end{enumerate}

\noindent
ほかに、教師が遅延に依らず下界（rms $<0.2$\,deg）であること、
\code{naive} が遅延とともに単調悪化すること、
齢の信号が実際に変動し輸送遅延以上であること、
文脈走査の O(1) 実装が総当たりと一致することを検査する。

\section{実験1：輸送遅延の掃引}

\code{DELAYS=0,50,100,200,400,800 SEEDS=5}。以下はすべて追従誤差 rms（deg）の中央値（5 種）。

\begin{table}[htbp]
\centering
\small
\caption{実験1の結果。rms（deg）の中央値}
\label{tab:e1}
\begin{tabular}{@{}r rr rr rr@{}}
\toprule
\textbf{遅延} & \textbf{teacher} & \textbf{naive} &
\textbf{P0 模倣・齢なし} & \textbf{P1 模倣・齢あり} &
\textbf{P2 RL・齢なし} & \textbf{P3 RL・齢あり} \\
\midrule
0\,ms   & 0.092 & 3.433 & 0.644 & \textbf{0.378} & 0.248 & 0.265 \\
50\,ms  & 0.092 & 4.894 & 0.945 & \textbf{0.769} & 0.795 & 0.727 \\
100\,ms & 0.092 & 6.325 & 1.595 & 1.436 & 1.725 & 1.858 \\
200\,ms & 0.092 & 9.045 & 3.286 & 3.059 & 3.114 & 3.233 \\
400\,ms & 0.092 & 13.727 & 7.333 & 8.929 & 6.634 & \textbf{6.456} \\
800\,ms & 0.092 & 19.100 & 15.938 & 15.860 & 12.381 & \textbf{13.533} \\
\bottomrule
\end{tabular}
\end{table}

\begin{figure}[htbp]
\centering
\begin{tikzpicture}
\begin{axis}[
  width=0.70\textwidth, height=66mm,
  xlabel={輸送遅延 $d$ [ms]}, ylabel={追従誤差 rms [deg]},
  ymode=log, log basis y=10,
  xmin=-30, xmax=830, ymin=0.07, ymax=30,
  grid=both, grid style={gray!22},
  legend style={at={(1.02,1)}, anchor=north west, font=\scriptsize, draw=gray!40},
  tick label style={font=\scriptsize}, label style={font=\small},
]
\addplot[thick,black,dashed,mark=none] coordinates
  {(0,0.0921)(50,0.0921)(100,0.0921)(200,0.0921)(400,0.0921)(800,0.0921)};
\addlegendentry{teacher（下界）}
\addplot[thick,gray,mark=square*,mark size=1.6] coordinates
  {(0,3.4331)(50,4.8943)(100,6.3253)(200,9.0453)(400,13.7274)(800,19.0998)};
\addlegendentry{naive}
\addplot[thick,orange!80!black,mark=o,mark size=1.8] coordinates
  {(0,0.6445)(50,0.9448)(100,1.5952)(200,3.2861)(400,7.3330)(800,15.9378)};
\addlegendentry{P0 模倣・齢なし}
\addplot[thick,red!70!black,mark=*,mark size=1.6] coordinates
  {(0,0.3778)(50,0.7693)(100,1.4362)(200,3.0589)(400,8.9294)(800,15.8600)};
\addlegendentry{P1 模倣・齢あり}
\addplot[thick,blue!45!black,mark=triangle,mark size=2.2] coordinates
  {(0,0.2476)(50,0.7949)(100,1.7250)(200,3.1141)(400,6.6339)(800,12.3811)};
\addlegendentry{P2 RL・齢なし}
\addplot[thick,blue!70!green,mark=triangle*,mark size=2.2] coordinates
  {(0,0.2648)(50,0.7268)(100,1.8577)(200,3.2331)(400,6.4562)(800,13.5327)};
\addlegendentry{P3 RL・齢あり}
\end{axis}
\end{tikzpicture}
\caption{遅延に対する追従誤差（対数軸）。学習はどの遅延でも \code{naive} を大きく下回るが、
teacher には遠い。400\,ms 以上で強化学習（青系）が模倣（赤系）を明確に下回る。}
\label{fig:e1}
\end{figure}

\subsection{対応のある比較}

\prob{中央値だけでは、ばらつきに埋もれた差を見分けられない。}
同じ種どうしで比べ、\keyp{勝った種の数}を併記する。

\begin{table}[htbp]
\centering
\small
\caption{同一種での比較。括弧内は5種のうち改善した数}
\label{tab:paired}
\begin{tabular}{@{}r lll@{}}
\toprule
\textbf{遅延} & \textbf{齢の効果・模倣 P0$\to$P1} & \textbf{齢の効果・RL P2$\to$P3} &
\textbf{RL 対 模倣 P1$\to$P3} \\
\midrule
0\,ms   & \keyp{$-41\,\%$（5/5）} & $+7\,\%$（0/5） & \keyp{$-30\,\%$（5/5）} \\
50\,ms  & \keyp{$-19\,\%$（5/5）} & $-9\,\%$（4/5） & $-6\,\%$（3/5） \\
100\,ms & $-10\,\%$（4/5） & $+8\,\%$（2/5） & $+29\,\%$（0/5） \\
200\,ms & $-7\,\%$（4/5） & $+4\,\%$（2/5） & $+6\,\%$（2/5） \\
400\,ms & \prob{$+22\,\%$（1/5）} & $-3\,\%$（2/5） & \keyp{$-28\,\%$（5/5）} \\
800\,ms & $-0.5\,\%$（3/5） & $+9\,\%$（0/5） & \keyp{$-15\,\%$（5/5）} \\
\bottomrule
\end{tabular}
\end{table}

\subsection{読み取り}

\begin{enumerate}[leftmargin=1.8em,itemsep=4pt]
  \item \keyp{学習は明確に効く。} どの遅延でも \code{naive} を大きく下回る
        （800\,ms で 19.1 $\to$ 12.4\,deg）。ただし teacher（0.092\,deg）には遠い。
  \item \keyp{齢は、模倣では小さい遅延で効く。} 0\,ms で $-41\,\%$（5/5）、
        50\,ms で $-19\,\%$（5/5）。遅延が伸びるほど効果は薄れ、
        \prob{400\,ms では逆に悪化する}（$+22\,\%$、勝ち 1/5）。
  \item \prob{強化学習では齢の効果が検出できない。} 勝ち数は 0/5〜4/5 に散らばり、
        中央値の変化も $\pm 9\,\%$ 以内で、種のばらつきに埋もれている。
  \item \keyp{大遅延では強化学習が模倣を明確に上回る。}
        400\,ms で $-28\,\%$（5/5）、800\,ms で $-15\,\%$（5/5）。
        これは案Eの中心仮説 ――\keyp{教師が失敗する領域は模倣では学べず、そこは強化学習が取る}――
        を支持する、唯一かつ明瞭な結果である。
\end{enumerate}

\section{事前に宣言した境界の扱い}

\begin{itemize}[leftmargin=1.4em,itemsep=3pt]
  \item $d^{*}_{IL}$ ＝ \keyp{200\,ms}（P1 の rms が \code{TOL}$=2.0$\,deg を超える最小の遅延）。
        宣言どおり機能した。
  \item \prob{$d^{*}_{RL}$ の定義は機能しなかった。}
        「P3 が P1 を下回る最小の点」は 0\,ms を返すが、100\,ms と 200\,ms では逆転しており
        \prob{単調でない}ので、「境界」という語に値しない。
\end{itemize}

\noindent
\keyp{定義を後から都合よく変えない。} ここでは事実を書くにとどめる ――
\keyp{強化学習の優位は 400\,ms 以上で 5/5 種と安定する}。
次の実験では「$k$ 種すべてで下回る最小の点」という形で\emph{宣言し直して}測る。

\section{この軸では仮説を検定できない（設計上の発見）}
\label{sec:flaw}

\prob{遅延 $d$ は1回の走行のあいだ一定である。}
したがって齢 $a$ は「$d$（定数）＋配布周期 100\,ms のなかの位相」であり、
\keyp{情報として効くのは位相の部分だけ}で、定数 $d$ は学習で重みに吸収される。
位相の変動幅は $d$ によらず 100\,ms で一定なので、
\prob{「遅延が大きいほど齢が効く」はこの軸では原理的に出ない}。

実際、齢の効果は $d$ とともに薄れ、400\,ms では
\keyp{入力がほぼ定数になったぶん、むしろ害になった}（$+22\,\%$）と読める。
入力を1本増やせば、その情報が使えないとき過学習の材料になるだけである。

\medskip
\noindent
案Eの主張は「\keyp{遅延が変動する条件下で}」であった。
したがって掃引すべき軸は遅延の\emph{大きさ}ではなく\keyp{遅延のばらつき}である。
実験1は\prob{主張と対応しない軸を掃引していた}ことになる。

\medskip
\noindent
これは負の結果ではなく、\keyp{測る前に気づけなかった設計の誤り}である。
記録しておく価値があるのは、案Dの計画書にも同じ穴があるからである ――
案Dは遅延を「制御変数として掃引する」と書いているが、
\keyp{掃引すべきは分布の位置ではなく形（分散）}である。
案Dの実験計画をこの点で修正する必要がある。

\section{実験2：遅延のばらつきを掃引する}
\label{sec:e2}

\code{DELAYS=100 JITTERS=0,25,50,100,200 SEEDS=5}（輸送遅延は 100\,ms に固定）。
ジッタは文脈の生成時刻に $\pm j$ の一様乱数として入る。
追い越して着いた古い文脈は\keyp{最新優先で捨てる}
（案Aの「文脈クラスは最新優先で上書き」に対応する方針を実装した）。
これにより齢の変動幅は $j=0$ の 100\,ms から $j=200$\,ms では \keyp{279\,ms} へ広がり、
\keyp{齢が定数でなくなる}。

\subsection{予測（結果を見る前に書いたもの）}

\begin{enumerate}[leftmargin=1.8em,itemsep=2pt]
  \item ジッタが増えるほど\keyp{P1 対 P0 の差が開く}。齢が定数でなくなるため。
  \item $\hat a$ の効果はジッタで弱まる。次の到着時刻の予測が外れるようになるため。
  \item 強化学習側（P3 対 P2）でも差が出るなら、齢の効果は学習法に依らないことになる。
        出ないなら、\prob{CEM の探索が齢を使う解に届いていない}疑いが残る。
\end{enumerate}

\subsection{$\hat a$ の実装を先に直す}
\label{sec:ahead}

結果を出す前に、\prob{$\hat a$ の作り方に誤りがあった}ので記しておく。
最初の実装は $\hat a$ を「\emph{次に本当に届く時刻}」から計算していた。
これは\prob{OS には出せない値}である。tickless ワンショットタイマが持っているのは
\keyp{名目の配布周期}であって、ジッタや脱落で実際の到着がずれることを OS は知らない。

\noindent
さらに、この実装には副作用があった。
最初の文脈が届く前（$t<d$）に「次の到着」として\prob{2 番目の事象}を返しており
（$0.181$\,s。正しくは 1 番目の $0.081$\,s）、
その結果 $\hat a$ が\prob{「まだ起動区間である」ことの目印}として漏れていた。
ずれは $t<0.1$\,s の 102 標本に限られるが、そこは文脈を1つも持たない高レバレッジな区間である。

\noindent
名目周期から作る実装に直したところ、$j=0$ での齢の効果は\keyp{$-10\,\%$ から $0\,\%$ へ消えた}。
\keyp{$\hat a$ を使わない P0 の値は 1.595 で完全に一致}し、P1 だけが $1.436 \to 1.592$ と動いたので、
差の出どころがこの1点であることは確かめられている。

\medskip
\noindent
\keyp{OS が何を答えられるかを厳密にしないと、齢の効果を過大評価する。}
これは実装の瑕疵であると同時に、案E の主張そのものに関わる教訓である ――
上位に見せてよいのは\emph{OS が保証できる値}だけであり、
シミュレータの中では「本当の次の到着」がいくらでも手に入ってしまう。
以下の数値はすべて\keyp{名目周期版（\code{AHEAD=nominal}）}である。
比較用に \code{AHEAD=oracle} も残してある。

\subsection{結果}

\begin{table}[htbp]
\centering
\small
\caption{実験2の結果。rms（deg）の中央値（5 種）。$d=100$\,ms 固定、$\hat a$ は名目周期版}
\label{tab:e2}
\setlength{\tabcolsep}{5pt}
\begin{tabular}{@{}r rr rr rr@{}}
\toprule
\textbf{ジッタ $j$} & \textbf{teacher} & \textbf{naive} &
\textbf{P0 模倣・齢なし} & \textbf{P1 模倣・齢あり} &
\textbf{P2 RL・齢なし} & \textbf{P3 RL・齢あり} \\
\midrule
0\,ms   & 0.092 & 6.325 & 1.595 & 1.592 & 1.725 & 1.858 \\
25\,ms  & 0.092 & 6.391 & 1.782 & 1.686 & 1.967 & 2.093 \\
50\,ms  & 0.092 & 6.561 & 2.083 & \textbf{1.931} & 2.358 & 2.353 \\
100\,ms & 0.092 & 6.944 & 2.942 & \textbf{2.350} & 3.316 & 3.098 \\
200\,ms & 0.092 & 7.778 & 4.198 & \textbf{2.899} & 3.612 & 4.290 \\
\bottomrule
\end{tabular}
\end{table}

\begin{figure}[htbp]
\centering
\begin{tikzpicture}
\begin{axis}[
  width=0.70\textwidth, height=62mm,
  xlabel={文脈生成のジッタ $j$ [$\pm$ms]}, ylabel={追従誤差 rms [deg]},
  xmin=-10, xmax=210, ymin=0, ymax=8.4,
  grid=both, grid style={gray!22},
  legend style={at={(1.02,1)}, anchor=north west, font=\scriptsize, draw=gray!40},
  tick label style={font=\scriptsize}, label style={font=\small},
]
\addplot[thick,gray,mark=square*,mark size=1.6] coordinates
  {(0,6.325)(25,6.391)(50,6.561)(100,6.944)(200,7.778)};
\addlegendentry{naive}
\addplot[thick,orange!80!black,mark=o,mark size=1.8] coordinates
  {(0,1.595)(25,1.782)(50,2.083)(100,2.942)(200,4.198)};
\addlegendentry{P0 模倣・齢なし}
\addplot[very thick,red!70!black,mark=*,mark size=1.6] coordinates
  {(0,1.592)(25,1.686)(50,1.931)(100,2.350)(200,2.899)};
\addlegendentry{P1 模倣・齢あり}
\addplot[thick,blue!45!black,mark=triangle,mark size=2.2] coordinates
  {(0,1.725)(25,1.967)(50,2.358)(100,3.316)(200,3.612)};
\addlegendentry{P2 RL・齢なし}
\addplot[thick,blue!70!green,mark=triangle*,mark size=2.2] coordinates
  {(0,1.858)(25,2.093)(50,2.353)(100,3.098)(200,4.290)};
\addlegendentry{P3 RL・齢あり}
\end{axis}
\end{tikzpicture}
\caption{遅延の\keyp{ばらつき}に対する追従誤差。模倣では齢あり（太い赤）と齢なし（橙）の
差がジッタとともに単調に開く。実験1（図\ref{fig:e1}）で薄れていったのと逆向きである。}
\label{fig:e2}
\end{figure}

\begin{table}[htbp]
\centering
\small
\caption{同一種での比較。齢の効果はジッタとともに\keyp{単調に強まる}}
\label{tab:paired2}
\begin{tabular}{@{}r l c l c@{}}
\toprule
\textbf{ジッタ} & \textbf{齢の効果・模倣 P0$\to$P1} & \textbf{勝ち} &
\textbf{齢の効果・RL P2$\to$P3} & \textbf{勝ち} \\
\midrule
0\,ms   & $-0.2\,\%$ & 2/5 & $+7.7\,\%$ & 2/5 \\
25\,ms  & $-5.4\,\%$  & 3/5 & $+6.4\,\%$ & 1/5 \\
50\,ms  & $-7.3\,\%$ & \keyp{5/5} & $-0.2\,\%$ & 4/5 \\
100\,ms & $-20.1\,\%$ & \keyp{5/5} & $-6.6\,\%$ & 3/5 \\
200\,ms & \keyp{$-30.9\,\%$} & \keyp{5/5} & $+18.8\,\%$ & 1/5 \\
\bottomrule
\end{tabular}
\end{table}

\subsection{読み取り}

\begin{enumerate}[leftmargin=1.8em,itemsep=4pt]
  \item \keyp{予測1は的中した。} 齢の効果は $0\,\% \to -31\,\%$ と\keyp{単調に強まり}、
        $j\ge 50$\,ms では\keyp{5 種すべてで}齢ありが勝つ。
        \prob{齢が効くのは遅延が大きいときではなく、遅延がばらつくときである。}
  \item 別の見方をすると、齢は\keyp{劣化を食い止める}。
        \code{naive} 比の改善率は、齢なしが $75\,\% \to 46\,\%$ と崩れていくのに対し、
        齢ありは $75\,\% \to 63\,\%$ にとどまる。
        ばらつきが増えるほど、齢を渡すことの価値が上がる。
  \item \prob{予測3は外れ、強化学習では今回も齢の効果が出なかった}（勝ち 1/5〜4/5、符号も一定しない）。
        実験1と合わせて、\keyp{CEM が齢を使う解に届いていない}疑いが濃い。
        121 次元の分布探索で、8 本ある入力のうち 2 本の使い方だけを見つけるのは、
        勾配を使わない探索には荷が重いという読みが自然である。
        別のアルゴリズムでの確認が要る。
  \item \keyp{ジッタ 0 では齢の効果が無い}（$-0.2\,\%$、2/5）。
        一方、実験1の $d=0$ では $-41\,\%$（5/5）だった。
        どちらも齢の変動は「配布周期 100\,ms の位相」だけなので、
        \prob{この不一致は未解明である}。輸送遅延そのものが、位相情報の使いみちを変えている
        可能性がある（$d=0$ なら外挿の horizon は最大 100\,ms、$d=100$\,ms なら 200\,ms）。
        第\ref{sec:next}節の課題に加える。
\end{enumerate}

\noindent
なお、この軸でも事前宣言どおり $d^{*}_{IL}$ は測れた ―― \keyp{ジッタ 100\,ms}
（P1 の rms が \code{TOL}$=2.0$\,deg を超える最小の点）。
$d^{*}_{RL}$ は掃引範囲では未到達である（P3 は P1 を一度も下回らない）。

\section{実験3：$a$ と $\hat a$ を分離する（アブレーション）}
\label{sec:e3}

実験1・2 は $a$ と $\hat a$ を同時に与えていたので、
\keyp{案E に固有の入力 $\hat a$ が単独でどれだけ効くか}が言えなかった。ここを埋める。
容量は4条件とも同一（端子は常に 8 本、切るのは値だけ）。
学習は\keyp{模倣のみ}である ―― 実験1・2 で CEM に齢の効果が出ないことが分かっているので、
ここに CEM を足しても $\hat a$ の情報価値は測れない。

\begin{table}[htbp]
\centering
\small
\caption{アブレーションの4条件}
\begin{tabular}{@{}l l l@{}}
\toprule
\textbf{条件} & \textbf{与える入力} & \textbf{対応} \\
\midrule
A0 & どちらも無し & 実験1・2 の P0 \\
A1 & \keyp{$a$ のみ} & 案D が渡すぶん \\
A2 & \keyp{$\hat a$ のみ} & 案E に固有。OS の tickless タイマだけが答えられる値 \\
A3 & 両方 & 実験1・2 の P1 \\
\bottomrule
\end{tabular}
\end{table}

\subsection{予測（結果を見る前に書いたもの）}

\begin{enumerate}[leftmargin=1.8em,itemsep=2pt]
  \item ジッタ軸では A1 が主役。A2 単独の効果は A1 より小さい。
  \item 脱落が増えると\keyp{A2 の効果が落ちる}。次の到着時刻の予告が外れるため。
        A1 は落ちない（齢は実測値なので脱落しても正しい）。
  \item A3 は A1 と A2 の和より小さい（情報が重複しているため）。
\end{enumerate}

\subsection{結果}

結果を表\ref{tab:e3}に示す。ジッタ軸（上4行）と脱落軸（下2行）を1つにまとめた。
脱落軸はジッタ $\pm100$\,ms を固定して掃引しているので、
両軸の交点（ジッタ $\pm100$、脱落 0）は同じ試行である。

\begin{table}[htbp]
\centering
\small
\caption{アブレーション。rms（deg）の中央値と、A0 を基準にした改善率（勝った種の数）。
$d=100$\,ms 固定。上4行はジッタ軸、下2行は脱落軸（ジッタ $\pm100$\,ms 固定）}
\label{tab:e3}
\begin{tabular}{@{}l r rl rl rl@{}}
\toprule
\textbf{条件} & \textbf{A0 なし} &
\multicolumn{2}{c}{\textbf{A1 \ $a$ のみ}} &
\multicolumn{2}{c}{\textbf{A2 \ $\hat a$ のみ}} &
\multicolumn{2}{c}{\textbf{A3 \ 両方}} \\
\midrule
ジッタ 0        & 1.595 & 1.394 & $-12.6\,\%$ (3/5) & 1.463 & $-8.3\,\%$ (3/5) & 1.592 & $-0.2\,\%$ (2/5) \\
ジッタ $\pm$50   & 2.083 & 1.878 & $-9.8\,\%$ (5/5)  & 2.270 & $+9.0\,\%$ (1/5) & 1.931 & $-7.3\,\%$ (5/5) \\
ジッタ $\pm$100  & 2.942 & 2.306 & $-21.6\,\%$ (5/5) & 2.741 & $-6.8\,\%$ (5/5) & 2.350 & $-20.1\,\%$ (5/5) \\
ジッタ $\pm$200  & 4.198 & 3.022 & \keyp{$-28.0\,\%$ (5/5)} & 3.936 & $-6.3\,\%$ (5/5) & 2.899 & \keyp{$-30.9\,\%$ (5/5)} \\
\midrule
脱落 10\,\%      & 3.341 & 2.393 & $-28.4\,\%$ (5/5) & 3.117 & $-6.7\,\%$ (2/5) & 2.237 & $-33.0\,\%$ (5/5) \\
脱落 30\,\%      & 5.976 & 3.277 & \keyp{$-45.2\,\%$ (5/5)} & 4.985 & $-16.6\,\%$ (5/5) & 3.047 & \keyp{$-49.0\,\%$ (5/5)} \\
\bottomrule
\end{tabular}
\end{table}

\subsection{読み取り}

\begin{enumerate}[leftmargin=1.8em,itemsep=4pt]
  \item \keyp{予測1は的中。主役は齢 $a$ である。}
        ジッタ $\pm200$\,ms で A1 が $-28.0\,\%$、A2 は $-6.3\,\%$。
        ただし\keyp{$\hat a$ 単独でも符号は一貫して正}で、$j\ge100$\,ms では 5/5 種が勝つ。
        \keyp{案Eに固有の入力は、単独でも効くが主役ではない}というのが正直な結論である。
  \item \prob{予測2は外れた。} 脱落が増えると A2 の効果は落ちるはずだったが、
        脱落 30\,\% で\keyp{$-16.6\,\%$（5/5）と逆に上がった}。
        名目周期は外れても\keyp{「基準」としては使える}という読みになる
        （いつ来るはずだったかを知っていれば、来なかったことが分かる）。
  \item \keyp{効果が最大になるのは脱落条件である。}
        脱落 30\,\% で A3 が $-49.0\,\%$（5/5）。
        これは\keyp{メールボックス溢れ}（\code{cc.c:532}）を模した条件であり、
        案Aが「暗黙の取りこぼしを明示的な鮮度契約に置き換える」と述べた、
        まさにその状況で齢がもっとも効いている。
  \item \keyp{予測3は的中。効果は劣加法的である。}
        脱落 30\,\% で $-45.2\,\%$ と $-16.6\,\%$ に対し、両方でも $-49.0\,\%$ にしかならない。
        $a$ と $\hat a$ は同じ「文脈がいつのものか」を別角度から見た量なので、当然ではある。
\end{enumerate}

\section{限界}

\begin{itemize}[leftmargin=1.4em,itemsep=3pt]
  \item \prob{1 関節・剛体・摩擦なし}。バックラッシュ、静止摩擦、サーボの応答遅れは入っていない。
        実機微調整の段階で必ずずれる（案Eの2年目の項目）。
  \item \prob{\code{VMAX}=180\,deg/s は推定値}でカタログ値ではない（\code{rl/dofbot\_rl.py} と同じ仮定）。
  \item \prob{強化学習は CEM のみ}。純 Python で回る規模に収めるための選択で、
        SAC 等との比較は未実施。したがって「強化学習では齢が効かない」は
        \keyp{CEM では検出できなかった}と読むべきである。
  \item \prob{$\hat a$ は名目周期からの計算}であり、実機の tickless タイマが持つ精度・
        誤差はまだ入っていない（第\ref{sec:ahead}節）。
        実機では「名目どおりに起きたか」自体がずれるので、この値の質は測り直しになる。
  \item \prob{ジッタ 0 での不一致が未解明}。実験1の $d=0$ では齢が $-41\,\%$ 効くのに、
        実験2の $j=0$（$d=100$\,ms）では効かない。どちらも齢の変動は位相のみである。
  \item 模倣と強化学習は学習量を揃えていない（別アルゴリズムなので原理的に揃わない）。
        厳密な対照は \keyp{P0 対 P1} と \keyp{P2 対 P3} の2組であり、
        P1 対 P3 は境界の定義にのみ使う。
  \item \prob{齢は理想的に正確}。実機では OS の計時そのものに誤差がある。
        外部計測（ロジック・アナライザ）で裏を取る設計は案Eの第3項目にある。
  \item 課題は\keyp{追従のみ}。「到達して止まる」課題は入っていない。
        \keyp{「待つ」が最適になるのはそちら}なので、$\hat a$ の効果は過小評価されている可能性が高い。
\end{itemize}

\section{次の一手}
\label{sec:next}

\begin{enumerate}[leftmargin=1.8em,itemsep=3pt]
  \item \keyp{ジッタ 0 での不一致の解明}（$d=0$ で $-41\,\%$、$d=100$\,ms で $0\,\%$）。
        齢の情報が「外挿の horizon」として使われているなら、
        $d$ と外挿距離を分けて与える実験で切り分けられる。\keyp{最優先}。
  \item \keyp{到達して止まる課題}の追加。「待つ」が最適解になる課題で $\hat a$ を測り直す。
        追従課題では $\hat a$ の効果は過小評価されているはずである。
  \item 強化学習を CEM 以外でも回し、「齢が効かない」が手法固有でないことを確認する。
        勾配を使う手法なら 2 本の入力の使い方を見つけられるか、という問いである。
  \item 脱落条件をさらに詰める。\keyp{効果が最大だったのは脱落 30\,\% の $-49\,\%$}であり、
        ここが案Aの鮮度契約と直接つながる。落とし方の方針（最新優先・取りこぼし禁止）を
        条件に加えて測る。
  \item 多関節（隣接関節の齢つき情報をメッシュ経由で受ける）へ拡大する。
  \item 固定小数点・量子化した方策で同じ実験を回し、案F・案Gの入力にする。
\end{enumerate}

\section{案D・案Eの計画へ戻すべき修正}

本実験の結果は、計画書の側を2箇所直すことを求める。

\begin{enumerate}[leftmargin=1.8em,itemsep=3pt]
  \item \prob{掃引すべきは遅延の大きさではなく、ばらつきである。}
        案D（\code{vla\_rt\_review.pdf}）は「遅延を 0.1\,ms から数百 ms まで変えた軌跡を生成する」
        と書いているが、\keyp{固定遅延をいくら振っても齢の効果は測れない}
        （第\ref{sec:flaw}節）。分布の\emph{位置}ではなく\emph{形}を掃引する設計に直す必要がある。
  \item \keyp{齢の効果の主戦場は大遅延ではなく、ばらつきと脱落である。}
        したがって案Dの合格ラインも、遅延の大きさではなくばらつきの関数として書くべきである。
        実験2の $-31\,\%$（ジッタ $\pm 200$\,ms、5/5 種）と
        実験3の $-49\,\%$（脱落 30\,\%、5/5 種）が、その最初の目安になる。
  \item \keyp{上位に見せる値は「OS が保証できる値」に限る。}
        第\ref{sec:ahead}節で見たとおり、シミュレータの中では
        「本当の次の到着」がいくらでも手に入ってしまい、齢の効果を過大評価する。
        案A の \code{ctx\_age\_us()} を設計するときも、
        \prob{OS が実際に答えられる値の範囲を先に決めてから}方策に渡すこと。
\end{enumerate}

\section*{付録：構成と再現手順}

\begin{table}[htbp]
\centering
\small
\caption{\code{xinu-mesh-rtos/planE/}（純 Python、依存ゼロ）}
\begin{tabularx}{\textwidth}{@{}l X@{}}
\toprule
\code{joint.py} & サーボ、上位目標、齢つき文脈配布（脱落・ジッタ・追い越し処理を含む） \\
\code{policy.py} & 121 パラメータの残差方策、齢マスク、Adam \\
\code{env.py} & 1 エピソードの実行と、分位点つきの指標 \\
\code{train.py} & 特権教師つき模倣（DAgger 1 巡）と CEM 強化学習 \\
\code{sweep.py} & 遅延・ジッタ・脱落率の掃引と要約。しきい値をここで宣言する \\
\code{test\_planE.py} & 実験系そのものの検証（17 項目） \\
\code{ablate.py} & $a$ と $\hat a$ を分離するアブレーション（実験3） \\
\code{plot\_data.py} & \code{results/*.csv} を pgfplots 用の表に落とす \\
\bottomrule
\end{tabularx}
\end{table}

\begin{verbatim}
python3 test_planE.py                                   # 先にこれを通す
DELAYS=0,50,100,200,400,800 SEEDS=5 python3 sweep.py
DELAYS=100 JITTERS=0,25,50,100,200 SEEDS=5 OUT=sweep_jitter python3 sweep.py
JITTERS=0,50,100,200 SEEDS=5 python3 ablate.py                 # 実験3（ジッタ軸）
PDROPS=0,0.1,0.3 JITTERS=100 SEEDS=5 OUT=ablate_drop python3 ablate.py
\end{verbatim}

\noindent
所要はいずれも約 35 分（Apple Silicon、単一プロセス）。
生データは \code{results/*.csv}、要約は \code{results/*\_summary.md} に出る。
指標は平均だけでなく p50/p90/p99/max を必ず出す ―― \keyp{集計値は嘘をつく}。

\end{document}
